% !PRINTED 517
\weberpage{517}
If $\omega$ is prime, the theorem is evident, because $\omega$ is divisible by itself. But if $\omega$ is not prime, then it is divisible by an integral functional $\omega_1$ which is neither a unit nor associated with $\omega$. If, then,
\weq{1}{\omega = \omega_1 \alpha,}
neither $\omega_1$ nor $\alpha$ is a unit. But from (1) follows
\[
N_a(\omega) = N_a(\omega_1)\, N_a(\beta), % [sic] printed with $\beta$, though (1) introduces $\alpha$
\]
and since $N_a(\beta)$ is greater than 1, we have
\weq{2}{N_a(\omega_1) < N_a(\omega).}

If $\omega_1$ is not yet prime, one can apply the same inference to $\omega_1$ and finds a divisor $\omega_2$ of $\omega_1$ such that
\weq{3}{N_a(\omega_2) < N_a(\omega_1) < N_a(\omega)}
holds. But since there are only a finite number of natural numbers that are smaller than $N_a(\omega)$, the sequence of functionals $\omega, \omega_1, \omega_2, \ldots$ must break off, and this is possible only if the last of them is a prime functional, whereby Theorem 1. is proved.

2. \gesperrt{Every integral functional $\omega$ in the field $\Omega$ that is different from zero and from the units can be decomposed into a finite number of prime factors.}

For if $\pi_1$ is a prime factor of $\omega$ and
\weq{4}{\omega = \pi_1 \omega_1,}
then, since $N_a(\pi_1) > 1$, we have
\[
N_a(\omega) > N_a(\omega_1).
\]

If $\omega_1$ is not a unit, it is divisible by a prime functional $\pi_2$, and from
\[
\omega_1 = \pi_2 \omega_2
\]
follows
\[
N_a(\omega_1) > N_a(\omega_2).
\]

Continuing in this way, one obtains a sequence of integral functionals $\omega_1, \omega_2, \ldots$, whose absolute norms continually decrease, and this sequence therefore breaks off with a unit. If $\pi_\nu$ is the last of them that is not a unit, then $\pi_\nu$ is itself a prime functional, and it follows that
\weq{5}{\omega = \pi_1 \pi_2 \ldots \pi_\nu,}
q. e. d.
